% ---------------------------------------------------------------------------------------------
% Main text, part 2 (Sections 4-5). Numbers from tables_tax.tex / _res_tax_*.csv (2026-10-03).
% ---------------------------------------------------------------------------------------------

\section{Results}
\label{sec:results}

\subsection{Seats: the periphery loses, the hubs do not}

Table~\ref{tab:seats} reports the first-year effect of a ticket tax on log scheduled seats. With equal
airport weights, seats at taxed airports fall by 10.3 percent relative to control airports (column 1,
$p<0.05$). With seat weights the effect is 4.1 percent (column 2, $p<0.10$). The difference is the
distribution of the effect across the size terciles in columns 4 to 6: airports in the smallest tercile of
the taxing country lose 24.0 percent of their seats, the middle tercile 1.6 percent (not significant)
and the largest tercile 4.8 percent (not significant). The pre-trend coefficient for the second
pre-announcement year is $-0.031$ with a standard error of 0.021, and the pooled event study in
Table~\ref{tab:eventtime} shows pre-period coefficients of $-0.051$ and $-0.031$ (both insignificant)
against post-period coefficients of $-0.124$ in the first tax year and $-0.153$ in the second (both
$p<0.01$). Event by event, the first-year coefficients are $-0.18$ for Germany, $-0.12$ for Austria,
$-0.13$ for Norway and $-0.16$ for Sweden, all significant, against $-0.01$ for the UK rate increase and
$-0.10$ (not significant) for the Netherlands; the UK pre-period shows a trend and the UK event carries
the smallest dose.

The border group does not gain. Foreign airports 50 to 300 kilometres from the taxed country show
coefficients of $+0.009$ (equal weights) and $-0.030$ (seat weights), neither consistent with the
inflow of passengers that leakage would produce. Table~\ref{tab:robust} shows that the equal-weighted
estimate lies between $-0.11$ and $-0.15$ across effective-date windows of 12 to 36 months, without the
two recession-year events, with announcement donuts of two or three years and after removing
airport-specific pre-trends.

\subsection{Network position: concentration, not contraction}

Table~\ref{tab:network} turns to the annual network outcomes. The connectivity index of taxed airports
falls by 1.6 percent (equal weights, $p<0.05$), 2.3 percent (seat weights, $p<0.01$) and 2.6 percent at
hubs with more than one million pre-year seats. The components show where the index moves: taxed
airports lose 8.9 percent of their destinations, 18.6 percent of eigenvector centrality and 42.7
percent of betweenness centrality. By size, the index falls by 2.3 percent in the small and mid terciles
and by an insignificant 0.3 percent in the large tercile; destinations fall by 13.5 and 10.6 percent in
the small and mid terciles against 3.4 percent (not significant) at large airports; eigenvector
centrality falls by 32 percent at small airports and 13 percent at mid and large airports. The annual
event study (Table~\ref{tab:eventtime}, Panel B) is the cleanest evidence in the paper: the GACI
coefficients are $-0.001$ (SE 0.008) three years and $0.000$ (SE 0.003) two years before the effective
year, $-0.016$ in the effective year and $-0.018$ the year after (both $p<0.01$). Border airports show
no change in connectivity.

Two features of these estimates matter for interpretation. First, GACI is standardised within each
year, so the fall is relative: taxed airports slipped in the world ranking while the ranking of untaxed
airports improved. Second, the losses in degree and eigenvector centrality at small and mid-sized
airports, with hubs holding their degree, describe a network that is thinning at the edges and
concentrating at the centre. The betweenness losses are large in proportional terms because the
betweenness of small airports is close to zero and volatile; the point estimate for large airports is
insignificant.

\subsection{Emissions fall where flights disappear}

Table~\ref{tab:emissions} reports the same specifications for emissions. Departure CO$_2$ at taxed
airports falls by 10.4 percent with equal weights ($p<0.05$) and by 2.1 percent with seat weights (SE
0.021, not significant). By size, small airports lose 21.2 percent of their emissions ($p<0.01$), mid
airports 3.1 percent and large airports 6.4 percent (both not significant). Emissions per seat do not
fall: the seat-weighted coefficient is $+0.021$ ($p<0.05$), and emissions per seat-kilometre rise by
1.3 percent seat-weighted ($p<0.10$) and by 5.2 percent at small airports ($p<0.05$). Stage length is
unchanged everywhere. The tax therefore removes flying without making the flying that remains cleaner;
if anything, the loss of thin routes raises the average emissions intensity of the survivors.

Table~\ref{tab:mechanism} shows how. Panel A applies the exact decomposition \eqref{eq:decomp} to the
annual stacks. The total CO$_2$ effect of $-0.100$ splits into $-0.097$ from destinations, $-0.010$
from flights per destination and $+0.007$ from CO$_2$ per flight; for small airports the split is
$-0.161$, $-0.052$ and $+0.027$ of a total of $-0.186$. The entire emissions reduction is route exit.
Panel B confirms the margins at monthly frequency: small airports lose 23.7 percent of their flights,
keep aircraft size unchanged and are 4.4 percentage points more likely to record a month with no
scheduled departure at all; large airports lose 7.3 percent of flights ($p<0.10$) and raise seats per
flight by 2.5 percent ($p<0.05$). A flat per-passenger tax raises the fare on a thin, low-fare route by
a larger fraction than on a trunk route, so airlines cut the thin routes; at hubs, where transfer
passengers are exempt and the tax is a small share of the fare, they consolidate frequencies onto
larger aircraft.

\subsection{Connectivity and emissions}

Table~\ref{tab:gaci_co2} links the two sets of results in the manner of Table~9 in
\citet{li2019subway}, which translates network-density changes into pollution changes with the
estimated elasticity. Panel A estimates the within-airport elasticity of emissions to connectivity: a
one percent increase in GACI is associated with a 4.7 percent increase in CO$_2$ with airport and year
fixed effects and 3.6 percent with airport and country-by-year fixed effects, a 0.5 percent increase in
seats per flight, a 0.7 percent increase in stage length and a 0.5 percent fall in emissions per
seat-kilometre: more connected airports emit more in total and less per seat-kilometre. Multiplying
the tax effect on GACI from Table~\ref{tab:network} ($-0.016$ to $-0.023$) by this elasticity implies a
CO$_2$ change of 6 to 11 percent, in line with the direct equal-weighted estimate of $-10.4$ percent in
Table~\ref{tab:emissions}. Panel B makes the same point airport by airport with the
\citet{gelbach2016} decomposition: for small airports the total effect of $-0.186$ splits into $-0.069$
through the change in GACI and $-0.116$ holding GACI fixed; for mid airports $-0.073$ runs through
GACI and the direct part is zero; for large airports the direct part ($-0.052$) dominates a small
network component. Emissions intensity shows no network component anywhere. Because the elasticity
in Panel A is not instrumented within the tax stacks, we read Panel B as a consistent accounting of the
channel rather than as a structural mediation estimate (Appendix).

\subsection{National totals barely move}

The airport-level estimates say that taxes redistribute flying away from the periphery. Whether they
reduce national flying is a different question, and Table~\ref{tab:sdid} answers it with synthetic
difference-in-differences on national totals. For seats, the estimates are $-0.005$ for Germany (placebo
$p=1.00$), $-0.038$ for Austria ($p=0.83$), $-0.009$ for Sweden ($p=1.00$), $-0.016$ for the UK
($p=0.80$) and $-0.076$ for the Netherlands ($p=0.31$); only Ireland ($-0.159$, $p<0.01$), Malta
($-0.207$, $p<0.01$) and Norway ($-0.155$, $p=0.07$) show national losses, the first two coinciding with
the deepest national recessions in the sample and the third with a high flat tax in a market dominated
by domestic flying. On the seat-weighted national connectivity index, no event moves more than 5
percent and none is significant at 5 percent. Taken with the airport results, the picture is that of a
tax that reshuffles seats and network position within the country toward its hubs while leaving the
national total, and the country's position in the world network, largely where they were.

\subsection{Leakage}

Table~\ref{tab:leakage} splits the border group by distance to the nearest treated airport. Foreign
airports 50 to 150 kilometres away show $-0.019$ (SE 0.042) and those 150 to 300 kilometres away
$-0.002$ (SE 0.030); the 0 to 50 kilometre ring, which holds few airports and shows pre-trends, is
$-0.202$ ($p<0.10$). Splitting the 150 kilometre ring by size gives $0.000$ for foreign hubs with more
than five million seats and $-0.037$ for secondary airports. The one case with measurable leakage is
the Netherlands: while the 2008 tax was in force, seats at foreign airports 50 to 150 kilometres from
the border rose by 16.5 percent ($p<0.01$), and 14.9 percent in the eighteen months after the tax was
set to zero, consistent with the shift of Dutch passengers to Weeze, D\"usseldorf and Brussels
documented by \citet{gordijn2011}. The Dutch case combines a high rate, a dense foreign airport system
within an hour's drive and a one-year horizon; the other eight events show no detectable displacement
across borders.

\subsection{Cost per tonne}

Table~\ref{tab:abatement} turns the first-year estimates into an accounting of abatement, revenue and
connectivity lost, in the spirit of the cost-benefit section of \citet{li2019subway}. Applying the size-specific CO$_2$ coefficients to pre-year emissions, the 72 small taxed airports avoided 42 kt of CO$_2$
in the first tax year while generating EUR 20 million of revenue (at a load factor of 0.8), or EUR 484
of revenue per tonne avoided; the 67 mid airports avoided 53 kt at EUR 2,598 per tonne; and the 70
large airports, whose coefficient is not statistically different from zero, account for 3,836 kt and EUR
972 per tonne at the point estimate. Over all treated airports the revenue per tonne is EUR 988, an
order of magnitude above the EU allowance price of EUR 5 to 25 in the years of these events and five to
ten times the social cost of carbon of EUR 100 to 200 used in current appraisal. The abatement is also
bought with connectivity: small airports lose 0.021 index units and one destination per kilotonne
avoided, large airports 0.00007 index units and 0.04 destinations. The cost of the abatement a ticket
tax delivers is paid in the connectivity of peripheral regions.

\section{Discussion and conclusion}
\label{sec:conclusion}

Using nine European ticket-tax events, a monthly panel of more than 6,000 airports, flight-stage
emissions and a network-centrality index of airport position, this paper asks whether air passenger
taxes shrink the air network or concentrate it. The answer is the second. In the first tax year, taxed
airports lose 10 percent of their seats and emissions with equal weights but 4 and 2 percent with seat
weights; the smallest third of airports in the taxing country lose a quarter of their seats, a seventh of
their destinations and a third of their eigenvector centrality, while the largest third lose nothing and
upgauge. National seat totals and national network position are unchanged in synthetic
difference-in-differences for six of nine events. Emissions fall only through route exit; emissions per
seat-kilometre do not fall and rise at the hubs. Leakage across borders is absent except in the Dutch
episode.

Three implications follow. First, the flat per-passenger design with a transfer exemption is a tax on
thin, short, low-fare routes. It is therefore a tax on peripheral connectivity rather than on carbon:
the routes it removes are the ones on which the tax is a large share of the fare, not the ones with
the highest emissions. A levy proportional to fuel or distance, or an allowance price on emissions,
would target the intensity margin that the ticket tax leaves untouched, as the EU ETS results of
\citet{fageda2022pricing} suggest, at the cost of the hub exposure that \citet{fageda2025hub}
document. Second, the revenue a ticket tax raises per tonne avoided is far above any social cost of
carbon, which is consistent with the taxes' fiscal origins; if the objective is emissions, the
instrument is poorly matched to it. Third, because the losses fall on small and mid-sized airports and
the gains in position accrue to hubs, the tax raises the concentration of the national network. Our
companion work finds that hub concentration is associated with wider regional disparities, so the
distributional cost of the tax is likely to be regional as well as sectoral.

The estimates have limits. They are first-year effects on scheduled capacity and schedule-based
emissions; passengers, load factors and the longer-run adjustment of airline networks are not observed.
Identification rests on 33 country clusters and eight to nine events, and we therefore report event-by-event estimates, placebo inference for the national totals and sensitivity to parallel-trend violations rather than a single pooled standard error. The connectivity index is a relative measure, so the network result is a statement about position, not about the size of the world network. Finally, the
connectivity channel in Section~4.4 is an accounting decomposition; a credible instrument for airport
connectivity within short tax windows would allow the indirect effect to be identified rather than
described. Within these limits, the results give a clear answer to the question in the title: ticket
taxes cut carbon at the edges of the network and leave the centre intact, and the carbon they cut is
paid for in the connectivity of the periphery.
